\section{Un échéancier lié aux taux}
\label{sec:U3}

Les deux unités précédentes supposaient un taux donné, unique et constant. On
lève ici cette hypothèse : le taux devient une \emph{donnée de marché}, il dépend
de l'échéance, et il bouge. L'échéancier, lui, ne change pas --- ce qui change
est ce avec quoi on l'actualise.

C'est le premier étage où l'on rencontre un objet coté : l'obligation. Elle sert
de laboratoire, parce qu'elle rend visible sur un actif ce qui restera invisible
sur un passif d'assurance.

\subsection{La notion : prix d'une obligation à taux fixe}

\begin{notion}{Une obligation est un échéancier}
Une obligation à taux fixe de nominal $N$, de taux facial $c$ et de maturité $n$
verse un coupon $C = c\,N$ en fin de chacune des $n$ périodes, puis rembourse le
nominal à l'échéance. Son prix est la valeur actuelle de cet échéancier :
\[
  P \;=\; \sum_{k=1}^{n} \frac{C}{(1+r)^{k}} \;+\; \frac{N}{(1+r)^{n}}
      \;=\; C \times a_{\overline{n}|r} \;+\; N\,(1+r)^{-n}.
\]
Aucune formule nouvelle : c'est la VAN de l'unité~2 appliquée à une rente
certaine augmentée d'un flux terminal.

\medskip
La position du prix par rapport au nominal se lit immédiatement sur le rapport
entre taux facial et taux de marché : si $c > r$, l'obligation cote
\emph{au-dessus du pair} ; si $c < r$, \emph{au-dessous}.
\end{notion}

\begin{center}
\begin{tikzpicture}[x=1.25cm, y=0.8cm]
  \draw[timeaxis] (-0.4,0) -- (5.4,0);
  \foreach \x/\l in {0/0, 1/1, 2/2, 3/3, 4/4, 5/n}
    { \node[timenode] at (\x,0) {}; \node[etiq, below=3pt] at (\x,0) {\l}; }
  \draw[fluxneg] (0,0) -- (0,-1.0);
  \node[etiq, below] at (0,-1.05) {$P$};
  \foreach \x in {1,2,3,4}
    { \draw[flux] (\x,0) -- (\x,0.75); \node[etiq, above] at (\x,0.75) {$C$}; }
  \draw[flux] (5,0) -- (5,1.7);
  \node[etiq, above] at (5,1.7) {$C+N$};
\end{tikzpicture}
\end{center}

\begin{exercice}[prix d'une obligation]
Une obligation de nominal $1\,000$~\euro{} porte un taux facial de $3{,}20\,\%$
et arrive à échéance dans $5$~ans. Le taux de marché est de $2{,}60\,\%$.
Calculer son prix et commenter sa position par rapport au pair.
\end{exercice}

\begin{corrige}
Le coupon annuel vaut $C = 32$~\euro. Avec $(1{,}026)^{-5} = 0{,}87956$ :
\[
  a_{\overline{5}|2{,}6\%} = \frac{1-0{,}87956}{0{,}026} = 4{,}63242,
  \qquad
  P = 32 \times 4{,}63242 + 1\,000 \times 0{,}87956 \approx 1\,027{,}79\ \text{\euro}.
\]
L'obligation cote au-dessus du pair, à $102{,}78\,\%$, parce que son taux facial
excède le taux de marché : elle sert un coupon plus généreux que ce que le
marché exige aujourd'hui, et cet avantage se paie.
\end{corrige}

\begin{exercice}[lecture inverse]
La même obligation se négocie à $985$~\euro. Déterminer son taux actuariel.
\end{exercice}

\begin{corrige}
On cherche $r$ tel que $32\,a_{\overline{5}|r} + 1\,000\,(1+r)^{-5} = 985$.
Deux essais encadrent la solution :
\[
  P(3{,}50\,\%) = 986{,}45\ \text{\euro},
  \qquad
  P(3{,}60\,\%) = 981{,}99\ \text{\euro}.
\]
Une interpolation linéaire donne
$r \approx 3{,}50 + 0{,}10 \times \tfrac{986{,}45-985}{986{,}45-981{,}99}
 \approx 3{,}53\,\%$.

C'est le taux de rendement interne de l'unité~2, appliqué à une obligation. On
l'appelle alors \emph{taux actuariel} ou \emph{rendement à l'échéance}.
\end{corrige}

\subsection{Généralisation : courbe des taux, coupon couru, sensibilité}

\begin{generalisation}{Trois compléments indispensables}
\textbf{La courbe des taux.} Il n'existe pas un taux, mais un taux par échéance.
Le taux \emph{zéro-coupon}\index{taux zéro-coupon} $z_k$ est celui qui actualise
un flux unique disponible en $k$. Le prix exact s'écrit alors
\[
  P \;=\; \sum_{k=1}^{n} \frac{C_k}{(1+z_k)^{k}} .
\]
Le taux actuariel $r$ n'est pas une donnée de marché : c'est le taux unique qui,
appliqué à tous les flux, reproduit ce prix. C'est un \emph{résumé} de la courbe,
pas un ingrédient du calcul.

\medskip
\textbf{Le coupon couru.} Entre deux détachements, le porteur acquiert
progressivement le droit au coupon suivant. Le \emph{coupon
couru}\index{coupon couru} est cette fraction acquise :
\[
  \mathrm{CC} \;=\; C \times \frac{\text{jours écoulés depuis le dernier détachement}}{\text{base}} .
\]
Les obligations se cotent \emph{pied de coupon}, c'est-à-dire hors coupon couru.
Le montant effectivement réglé est la somme des deux.

\medskip
\textbf{Duration et sensibilité.} La \emph{duration de
Macaulay}\index{duration} est la moyenne des dates de flux, pondérée par leurs
valeurs actuelles :
\[
  D \;=\; \frac{1}{P}\sum_{k=1}^{n} k \,\frac{C_k}{(1+r)^{k}} .
\]
Elle s'exprime en années et mesure la maturité \emph{économique} du titre. La
\emph{sensibilité}\index{sensibilité} $S = D/(1+r)$ en donne la lecture
opérationnelle : pour une variation $\Delta r$ du taux,
\[
  \frac{\Delta P}{P} \;\approx\; -\,S \times \Delta r .
\]
Cette relation est une approximation du premier ordre. La relation prix-taux
étant convexe, elle surestime toujours la baisse et sous-estime toujours la
hausse.
\end{generalisation}

\begin{center}
\begin{tikzpicture}[x=1.15cm, y=1.5cm]
  \draw[timeaxis] (0,0) -- (5.7,0);
  \draw[timeaxis] (0,0) -- (0,1.5);
  \node[etiq, right] at (5.4,-0.25) {échéance};
  \draw[thick, bleuactu] plot[smooth] coordinates
    {(1,0.55) (2,0.72) (3,0.88) (4,1.00) (5,1.10)};
  \foreach \x/\y/\l in {1/0.55/{2{,}0}, 2/0.72/{2{,}3}, 3/0.88/{2{,}6},
                        4/1.00/{2{,}8}, 5/1.10/{3{,}0}}
    { \node[timenode] at (\x,\y) {}; \node[etiq, above left=-1pt] at (\x,\y) {\l}; }
  \foreach \x in {1,2,3,4,5} { \node[etiq, below=3pt] at (\x,0) {\x}; }
  \node[etiq, above right] at (0,1.5) {$z_k$ en \%};
\end{tikzpicture}
\end{center}

\begin{exercice}[courbe non plate et coupon couru]
\textbf{(a)} Reprendre l'obligation précédente --- nominal $1\,000$~\euro, taux
facial $3{,}20\,\%$, maturité $5$~ans --- et la valoriser avec la courbe
zéro-coupon ci-dessus : $2{,}0\,\%$, $2{,}3\,\%$, $2{,}6\,\%$, $2{,}8\,\%$ et
$3{,}0\,\%$.

\medskip
\textbf{(b)} L'obligation cote $101{,}05\,\%$ pied de coupon. L'achat intervient
$73$~jours après le dernier détachement, en base exact/365. Déterminer le montant
réglé par l'acheteur.
\end{exercice}

\begin{corrige}
\textbf{(a)} Chaque flux est actualisé à \emph{son} taux :
\begin{center}\small
\begin{tabular}{ccrr}
\toprule
Année & Flux & $z_k$ & Flux actualisé \\
\midrule
1 & $32$    & $2{,}0\,\%$ & $31{,}37$ \\
2 & $32$    & $2{,}3\,\%$ & $30{,}58$ \\
3 & $32$    & $2{,}6\,\%$ & $29{,}63$ \\
4 & $32$    & $2{,}8\,\%$ & $28{,}65$ \\
5 & $1\,032$ & $3{,}0\,\%$ & $890{,}21$ \\
\midrule
\multicolumn{3}{r}{\textbf{Prix}} & $\mathbf{1\,010{,}44}$ \\
\bottomrule
\end{tabular}
\end{center}
Le prix tombe à $1\,010{,}44$~\euro, contre $1\,027{,}79$~\euro{} avec un taux
plat de $2{,}6\,\%$. La pente de la courbe pénalise les flux longs, et le
remboursement du nominal représente ici près de $88\,\%$ du prix.

\medskip
\textbf{(b)} Coupon couru : $\mathrm{CC} = 32 \times \tfrac{73}{365} = 6{,}40$~\euro.
\[
  \text{montant réglé} = 1\,010{,}50 + 6{,}40 = 1\,016{,}90\ \text{\euro}.
\]
La cotation pied de coupon existe pour que le prix affiché ne chute pas
artificiellement à chaque détachement : elle isole la variation de valeur
économique du simple écoulement du temps.
\end{corrige}

\begin{exercice}[choc de taux, approché puis exact]
Reprendre l'obligation valorisée à $1\,027{,}79$~\euro{} au taux plat de
$2{,}60\,\%$.

\textbf{(a)} Calculer sa duration de Macaulay et sa sensibilité.

\textbf{(b)} Estimer l'effet d'une hausse instantanée des taux de $100$~points de
base par la formule de sensibilité, puis recalculer le prix exact. Commenter
l'écart.
\end{exercice}

\begin{corrige}
\textbf{(a)} On pondère chaque date par la valeur actuelle de son flux :
\begin{center}\small
\begin{tabular}{crrr}
\toprule
$k$ & Flux & Valeur actuelle & $k \times$ VA \\
\midrule
1 & $32$    & $31{,}19$  & $31{,}19$ \\
2 & $32$    & $30{,}40$  & $60{,}80$ \\
3 & $32$    & $29{,}63$  & $88{,}89$ \\
4 & $32$    & $28{,}88$  & $115{,}51$ \\
5 & $1\,032$ & $907{,}70$ & $4\,538{,}50$ \\
\midrule
  & & $1\,027{,}79$ & $4\,834{,}89$ \\
\bottomrule
\end{tabular}
\end{center}
\[
  D = \frac{4\,834{,}89}{1\,027{,}79} = 4{,}704\ \text{ans},
  \qquad
  S = \frac{4{,}704}{1{,}026} = 4{,}585 .
\]
La duration est inférieure à la maturité de cinq ans, puisque les coupons
avancent le centre de gravité des flux.

\medskip
\textbf{(b)} Estimation par la sensibilité :
\[
  \frac{\Delta P}{P} \approx -4{,}585 \times 0{,}01 = -4{,}585\,\%
  \quad\Longrightarrow\quad
  \Delta P \approx -47{,}12\ \text{\euro},
  \quad P \approx 980{,}67\ \text{\euro}.
\]
Calcul exact à $3{,}60\,\%$ : $P = 981{,}99$~\euro, soit $\Delta P = -45{,}80$~\euro,
c'est-à-dire $-4{,}457\,\%$.

\medskip
L'approximation surestime la perte de $1{,}32$~\euro, soit près de $3\,\%$ de
l'effet estimé. C'est la \emph{convexité} : la courbe prix-taux étant convexe,
la tangente passe sous elle des deux côtés. Sur un portefeuille obligataire de
plusieurs milliards et pour un choc réglementaire pouvant dépasser $100$~points
de base, cet écart cesse d'être une subtilité.
\end{corrige}

\begin{pont}{Où cela sert dans Solvabilité~II}
Cette unité se retrouve dans un bilan prudentiel à trois endroits distincts.

\medskip
\emph{À l'actif.} Les placements obligataires sont inscrits en juste valeur.
Le prix calculé ici est celui qui figure au bilan.

\emph{Au passif.} La courbe zéro-coupon de cette unité a une version
réglementaire : la structure par terme des taux sans risque publiée par l'EIOPA.
C'est elle, et non un taux plat, qui actualise les flux de prestations dans la
meilleure estimation.

\emph{Dans le capital requis.} Le sous-module de risque de taux consiste à
déformer la courbe à la hausse puis à la baisse et à recalculer le bilan complet.
L'exercice~3.4 en est la version élémentaire, appliquée à un seul titre.

\medskip
Enfin, la duration est l'instrument qui permet de comparer un actif et un passif
sur la même échelle. Une compagnie dont les engagements ont une duration de
douze ans et les placements une duration de cinq ans subit une perte lorsque les
taux baissent, quand bien même chaque poste, pris isolément, se comporterait
normalement. Cette confrontation est l'objet de la gestion actif-passif, traitée
à l'unité~7.
\end{pont}

\begin{refreg}
\textbf{Références.} Directive 2009/138/CE, article~75 : les actifs sont
valorisés au montant pour lequel ils pourraient être échangés dans le cadre
d'une transaction conclue entre parties informées et consentantes. \\
Article~77, paragraphe~2 : actualisation des flux de trésorerie futurs au moyen
de la courbe des taux sans risque pertinente, publiée mensuellement par l'EIOPA
avec sa documentation technique. \\
Règlement délégué (UE) 2015/35 : sous-module de risque de taux d'intérêt.
\emph{[Références d'articles à préciser avant diffusion.]}
\end{refreg}

\subsection*{Notations introduites}
\begin{description}[leftmargin=3cm, style=nextline, font=\normalfont]
  \notation{N}{nominal de l'obligation}
  \notation{c,\ C = cN}{taux facial et coupon}
  \notation{P}{prix de l'obligation}
  \notation{z_k}{taux zéro-coupon d'échéance $k$}
  \notation{\mathrm{CC}}{coupon couru}
  \notation{D}{duration de Macaulay, en années}
  \notation{S = D/(1+r)}{sensibilité}
\end{description}
